\section{대수적 대상(algebraic object): 몫환(quotient ring)에서 서명 바이트까지}
\label{sec:algebraic-object}

\subsection{기본 환(ring)과 mode~2 매개변수}

고정된 명세에서
\[
  R=\ZZ[X]/(X^{256}+1),\qquad
  R_q=\ZZ_q[X]/(X^{256}+1),\qquad q=64513
\]
이고, mode~2의 모듈 차원(module dimension)은 \((k,\ell)=(2,4)\), 작은
비밀계수(secret coefficient) 범위는
\(\eta=1\), binary challenge의 Hamming weight는 \(\tau=58\)이다
\cite{haetae-spec}. 각 \(R_q\) 원소는 차수 \(<256\)인 다항식(polynomial)의
동치류(equivalence class)이며, 계수벡터(coefficient vector)를 택하면
\(\ZZ_q^{256}\)와 집합(set)으로 식별할 수 있다.

명세는 \(q\equiv1\pmod{512}\)를 택해 \(X^{256}+1\)이 \(\FF_q\)에서 완전히
분해되고 NTT를 사용할 수 있게 한다. 따라서 종이 위에서는
중국인의 나머지 정리(Chinese remainder theorem)가 주는 동형사상(isomorphism)
\[
  R_q\simeq\prod_{i=1}^{256}\FF_q
\]
를 떠올릴 수 있다. 다만 현재 저장소의 최상위 구현 보안 인터페이스에는
\identifier{obl_public_key_algebra_ntt}가 별도 의무로 남아 있다. 즉 이
동형사상(isomorphism)의 존재와 실제 NTT 루틴의 정확성을 같은 주장으로
취급하면 안 된다.

\begin{caution}[수학 배경과 구현 정리(implementation theorem)의 분리]
이 절의 환(ring)·모듈(module) 서술은 고정된 HAETAE 명세를 설명한다. ``실제
Jasmin NTT가 이 동형사상(isomorphism)을 정확히 계산한다''는
구현 정리(implementation theorem)는
별도의 명제(proposition)와 전제(premise)가 필요하다. 본 안내서의
\statusproved 표시는 오직 \sourcepath{easycrypt/}
아래 명명된 정리(theorem)에 붙인다.
\end{caution}

\subsection{KeyGen을 모듈 사상(module map)으로 보기}

논문의 KeyGen 데이터 흐름에서 seed로부터
\(a_{\mathrm{gen}}\in R_q^k\)와
\(A_{\mathrm{gen}}\in R_q^{k\times(\ell-1)}\)를 만들고,
작은 원소(element)
\(s_{\mathrm{gen}}\in R^{\ell-1}\),
\(e_{\mathrm{gen}}\in R^k\)를 표본추출(sample)하여
\[
  b=a_{\mathrm{gen}}+A_{\mathrm{gen}}s_{\mathrm{gen}}
      +e_{\mathrm{gen}}\in R_q^k
\]
를 계산한다. 그 뒤 계수별 분해(coefficientwise decomposition)
\[
  b=b_0+2b_1,\qquad
  (b_0,b_1)=
  (\operatorname{LowBits}_{vk}(b),
   \operatorname{HighBits}_{vk}(b))
\]
를 사용하고, 검증키는 \(vk=(\mathit{seedA},b_1)\)로 압축된다. 비밀키의
packed representation은 Sign이 같은 검증키 성분을 다시 읽을 수 있도록
그 바이트 표현을 앞부분에 포함한다.

이 저장소에서 가장 먼저 닫힌 합성 seam은 전체 KeyGen 분포(distribution)가
아니라 바로 이
포함관계(inclusion relation)다. 바이트 배열의 집합(set)을 \(S\), 검증키
배열의 집합을 \(V\)라 하고
\[
  \pi_{992}:S\longrightarrow \Bytes^{992},\qquad
  \pi_{992}(s)=s[0..992)
\]
로 두면, 종료한 실제 mode~2 internal KeyGen 결과 \((vk,sk)\)에 대해
\[
  \pi_{992}(sk)=\pi_{992}(vk)
\]
가 성립한다. 오른쪽 배열은 더 크지만 첫 992바이트만 비교한다.

이를 몫(quotient)으로 정확히 표현할 수도 있다. \(S\)에
\[
  s\sim_{992}s'\quad\Longleftrightarrow\quad
  s[0..992)=s'[0..992)
\]
를 정의하면 \(\pi_{992}\)는 \(S/{\sim_{992}}\)를 분류한다. 현재
정리(theorem)는 ``packed secret key의 \(\sim_{992}\)-동치류(equivalence
class)가 public key의 표현과 같다''는 명제(proposition)다. 비밀키 전체의
등식(equality)이 아니다.

\subsection{Sign과 Verify가 만나는 transcript}

고정 명세에서 Sign과 Verify가 공통으로 재구성해야 하는 첫 관측값은 메시지
digest \(\mu\)다. 필요한 부분만 추리면
\[
  \mu=H_{\mathrm{gen}}(\mathit{seedA},b_1,M),
\]
그리고 challenge 입력은
\[
  \rho=H\bigl(
      \operatorname{HighBits}_h(w),
      \operatorname{LSB}(\lfloor y_0\rceil),
      \mu\bigr)
\]
의 모양을 갖는다 \cite{haetae-spec}. Sign은 packed secret key의 992바이트
접두부에서 \((\mathit{seedA},b_1)\)를 읽고, Verify는 public key 메모리에서
같은 바이트를 읽는다.

현재 raw/internal generated 경계에서 증명된 중심 등식은
\[
  \take_{32}\bigl(\mu_{\mathrm{Sign}}^{64}\bigr)
  =\mu_{\mathrm{Verify}}^{32}
\]
이다. 이어서 동일한 challenge 상태와 position~64에서 이 32바이트를 흡수하면
두 최종 상태와 position이 같다. 이 등식은 실제 generated hash/absorb 절차를
직접 포함하는 관계적 정리(relational theorem)다.

아직 증명되지 않은 것은 \(\rho\) 전체다. 특히
\(\operatorname{HighBits}_h(w)\)와 LSB 표현의 실제 Sign--Verify 일치,
전체 challenge 호출, 분포·rejection 분석은 별도 의무다. 그러므로 닫힌
\(\mu\) 사각형을 전체 challenge 또는 전체 \(\delta_{\mathrm{Sign}}=0\)로
확장할 수 없다.

\subsection{서명 표현은 정준 부분집합(canonical subset) 위의 부분동형(partial isomorphism)이다}

mode~2 서명은 1474바이트이고, 실제 추출된 layout은
\[
\underbrace{[0,32)}_{256\text{ challenge bits}}
\;\Vert\;
\underbrace{[32,1056)}_{1024\text{ signed low bytes}}
\;\Vert\;
\underbrace{[1056,1058)}_{\text{두 size delta}}
\;\Vert\;
\underbrace{[1058,1474)}_{\text{HBZ/h payload와 padding}}
\]
이다. 모든 bit word가 \(0\) 또는 \(1\)이고 모든 low word가 signed-byte
범위에 있는 정준 입력(canonical input)의 집합(set)을
\(C_{\mathrm{prefix}}\)라 하자. 현재 실제
prefix packer와 unpacker에 대해
\[
  \Decode_{\mathrm{prefix}}\circ
  \Encode_{\mathrm{prefix}}
  =\operatorname{id}_{C_{\mathrm{prefix}}}
\]
가 부분정확성(partial correctness)으로 증명되었다.

여기서 정의역 제한(domain restriction)은 장식이 아니다. noncanonical word에
대해서는 encode가 정보를 버리거나 decode가 다른 대표자(representative)를
돌려줄 수 있다. 대수학의 언어로는 \(C_{\mathrm{prefix}}\)가 선택된
대표자계(system of representatives)이고, round trip은 그 대표자계 위의
단면--수축 관계(section--retraction relation)다.

전체 서명 공간(signature space)에 대해서는 아직 이 등식(equality)이 없다.
size metadata, HBZ와 \(h\)
payload, canonical zero padding, 전체 parser의 성공과 유일성이 남아 있다.

\subsection{어떤 대수적 사실이 어느 층에 있는가}

\begin{center}
\begin{tabularx}{\textwidth}{L{31mm} Y L{30mm}}
\toprule
대상 & 필요한 명제(proposition) & 현재 상태 \\
\midrule
\(R_q\)-모듈(module) KeyGen & Jasmin 연산과 논문 KeyGen
분포(distribution)의 동일성 &
\statuspartial / 최상위 목표 \\
packed key & \(\pi_{992}(sk)=\pi_{992}(vk)\) &
\statusproved\ (partial correctness) \\
raw \(\mu\) transcript & Sign 64-byte 출력의 첫 32바이트와 Verify 32바이트 일치 &
\statusproved\ (generated raw/internal) \\
challenge suffix & 같은 position~64에서 \(\mu_{32}\) 흡수 뒤 상태 일치 &
\statusproved\ (국소 zero-loss) \\
signature prefix & canonical 입력에서 decode-after-encode가 항등(identity) &
\statusproved\ (prefix만) \\
HBZ/rANS suffix & 실제 full encoder와 decoder의 역관계(inverse relation) &
\statuspartial\ (actual core inverse는 \statusproved, production full wrapper는 미완) \\
전체 서명 스킴 & KeyGen/Sign/Verify refinement와 보안 합성 &
\statusspecified / \statusblocked \\
\bottomrule
\end{tabularx}
\end{center}

다음 절은 이 표의 ``부분정확성'', ``국소'', ``generated'', ``실제''가 어떤
논리적 차이를 갖는지 설명한다.
